\section{Numerical results}\label{sec:r24results}
The study contains 110 trajectories and 286 prescribed checkpoint records, separated into tuning, evaluation, multicell construction, an analytic reference and post-hoc diagnostics. These are not pooled as independent replications of a single claim. The supplement and machine-readable ledgers retain every configuration, candidate, termination and gate decision.

\subsection{Calibrated performance and the direct frontier}
Table~\ref{tab:r24tuning} charges every tuning trial. Every arm selects $.08$, the upper edge of the finite grid. Table~\ref{tab:r24heldout} reports final deployed policies, not favorable endpoints selected after evaluation. Equal calls do not imply equal work.
\begin{table}[htbp]\centering\small
\caption{Disjoint tuning, with all trials charged}\label{tab:r24tuning}
\begin{tabular}{lrrrrrr}\toprule
Arm & Rate & Calls & Gen. s & Prep. s & Check s & Mean lower payoff\\\midrule
Neural--Adam & 0.08 & 1200 & 3.944 & 0.000 & 12.420 & -1.287225\\
Direct--Adam & 0.08 & 1200 & 3.555 & 0.000 & 12.387 & -1.287524\\
Diagonal--Adam & 0.08 & 1200 & 3.570 & 0.018 & 12.474 & -1.287466\\
Tangent--Adam & 0.08 & 1200 & 3.547 & 0.018 & 12.425 & -1.287377\\
Whitened--Adam & 0.08 & 1200 & 3.535 & 0.018 & 12.375 & -1.304690\\
\bottomrule\end{tabular}
\par\smallskip\begin{minipage}{.97\textwidth}\footnotesize Two seeds, three rates, 200 calls per trial. Costs include every trial, not only the selected rate. Every arm selected the upper grid endpoint .08; optimal calibration beyond this declared finite grid is not established. Shared preprocessing is charged conservatively as a standalone cost.\end{minipage}
\end{table}

\begin{table}[htbp]\centering\small
\caption{Held-out central state: accuracy and computation}\label{tab:r24heldout}
\begin{tabular}{llrrrrr}\toprule
Rule & Arm & Regret upper & Calls & Gen. s & Check s & Cert. cost s\\\midrule
8/16 & Neural--Adam & 0.001925713 & 400.0 & 1.303 & 10.379 & 13.744\\
8/16 & Direct--Adam & 0.002039228 & 400.0 & 1.181 & 10.413 & 13.656\\
8/16 & Diagonal--Adam & 0.002008735 & 400.0 & 1.179 & 10.330 & 13.574\\
8/16 & Tangent--Adam & 0.002157232 & 400.0 & 1.175 & 10.333 & 13.573\\
8/16 & Whitened--Adam & 0.01874203 & 400.0 & 1.175 & 10.377 & 13.617\\
8/16 & Neural--L-BFGS-B & 0.009608327 & 251.0 & 0.792 & 10.478 & 13.333\\
8/16 & Direct--L-BFGS-B & 0.001923443 & 111.5 & 0.321 & 10.400 & 12.784\\
12/24 & Neural--Adam & 0.001925705 & 400.0 & 1.997 & 10.370 & 14.427\\
12/24 & Direct--Adam & 0.002039189 & 400.0 & 1.859 & 10.352 & 14.272\\
12/24 & Diagonal--Adam & 0.002008701 & 400.0 & 1.853 & 10.338 & 14.253\\
12/24 & Tangent--Adam & 0.002157182 & 400.0 & 1.847 & 10.377 & 14.286\\
12/24 & Whitened--Adam & 0.01874227 & 400.0 & 1.878 & 10.480 & 14.420\\
12/24 & Neural--L-BFGS-B & 0.009608181 & 243.5 & 1.183 & 10.479 & 13.722\\
12/24 & Direct--L-BFGS-B & 0.001923667 & 112.0 & 0.520 & 10.457 & 13.037\\
\bottomrule\end{tabular}
\par\smallskip\begin{minipage}{.97\textwidth}\footnotesize Regret is the worse of two seeds at (t,u,x)=(0,2,1.25); calls and times are means. Policies are final deployed checkpoints. Checking includes all five scheduled checks, including converged repeats. Certified cost adds preprocessing and initial checking. Diagnostic-only calls and tuning are separate.\end{minipage}
\end{table}

At the coarse rule, neural Adam's worst bound is $.001925713$, versus $.002039228$ for direct Adam and $.002008735$ for diagonal Adam. The frozen tangent chart does not reproduce the neural endpoint, and inverse-metric whitening is worse. These findings are consistent with roles for adaptive geometry, transported momentum and nonlinear steps, but do not separately identify those causal contributions.

Direct L-BFGS-B remains the strongest tested solver: its deployed bound is below $.001923443$ with mean $111.5$ calls and $.321$ generation seconds, versus $400$ calls and $1.303$ seconds for neural Adam. Neural L-BFGS-B retains an adverse held-out endpoint near $.009609$. The dual-free intervals in Table~\ref{tab:r24pairwise} identify payoff ordering without confounding it with a common unrestricted upper. The old shared-rate comparison remains historical; the present finite tuning envelope governs the new interpretation.
\begin{table}[htbp]\centering\small
\caption{Neural Adam minus comparator: dual-free payoff intervals}\label{tab:r24pairwise}
\begin{tabular}{llrr}\toprule
Rule & Comparator & Lower & Upper\\\midrule
8/16 & Direct--Adam & 0.0001125274 & 0.0001138956\\
8/16 & Diagonal--Adam & 0.00007642732 & 0.00008377674\\
8/16 & Tangent--Adam & 0.0001587773 & 0.0002319002\\
8/16 & Whitened--Adam & 0.01560812 & 0.01681708\\
8/16 & Neural--L-BFGS-B & $-2.561433\times10^{-6}$ & 0.007683369\\
8/16 & Direct--L-BFGS-B & $-2.798121\times10^{-6}$ & $-1.515012\times10^{-6}$\\
12/24 & Direct--Adam & 0.0001125009 & 0.0001138648\\
12/24 & Diagonal--Adam & 0.00007639466 & 0.00008375511\\
12/24 & Tangent--Adam & 0.0001587550 & 0.0002318586\\
12/24 & Whitened--Adam & 0.01560848 & 0.01681733\\
12/24 & Neural--L-BFGS-B & $-2.572662\times10^{-6}$ & 0.007683235\\
12/24 & Direct--L-BFGS-B & $-2.419365\times10^{-6}$ & $-1.474925\times10^{-6}$\\
\bottomrule\end{tabular}
\par\smallskip\begin{minipage}{.97\textwidth}\footnotesize Each row encloses both evaluation-seed payoff differences using independent delivered-policy intervals. Positive lower endpoints prove a payoff advantage in both matched problems; negative upper endpoints prove the reverse. This is not a confidence interval or an unrestricted optimality comparison.\end{minipage}
\end{table}


\subsection{Quadrature, geometry and the acceptance rule}
The largest ordinary raw-gradient difference between the two rules is \RXXIVGradientDiff{}, and the largest financing-offset difference is \RXXIVOffsetDiff{}. The largest relative difference is \RXXIVRelativeGradientDiff{}, at a nearly stationary fine-rule direct L-BFGS-B checkpoint. Small absolute differences therefore do not justify interpreting a surrogate stopping tolerance as a continuous-gradient tolerance. The principal endpoint ordering survives this finite refinement comparison, but no uniform stopped-gradient error bound follows. No held-out checkpoint node was clipped. Actual Jacobian and adaptive-metric spectra and transport remainders are retained as ordinary diagnostics, not rigorous rank certificates.

The 140 held-out replayed decisions contain 121 accepted replacements and 19 rejections: seven repeat the preceding candidate exactly and twelve concern changed candidates. The optimizer was not rolled back after rejection. These observations establish the recorded deployment effect, not a contribution to proposal convergence. Curves retain every scheduled checkpoint, including carried-forward converged endpoints, against actual calls, generation time and reconstructed generation-plus-checking cost.

\subsection{Expanded state coverage and known-solution calibration}
Table~\ref{tab:r24multi} reports all four cells through their union bound. Nine distinct policies per method are trained and checked. The neural union bound $.004280526$ and direct L-BFGS-B bound $.004274863$ both meet $.01$ throughout the expanded initialization region under the augmented-state realization. They are not sampled-point, high-dimensional or all-start-time results.
\begin{table}[htbp]\centering\small
\caption{Four cells and nine distinct experts per method}\label{tab:r24multi}
\begin{tabular}{lrrrr}\toprule
Arm & Region regret upper & Calls & Gen. s & Check s\\\midrule
Neural--Adam & 0.004280526 & 3600 & 11.755 & 18.686\\
Direct--Adam & 0.004398790 & 3600 & 10.553 & 18.638\\
Tangent--Adam & 0.004508256 & 3600 & 10.548 & 18.659\\
Direct--L-BFGS-B & 0.004274863 & 1250 & 3.592 & 18.694\\
\bottomrule\end{tabular}
\par\smallskip\begin{minipage}{.97\textwidth}\footnotesize Region [1.98,2.02] x [1.24,1.27] at time zero. Every cell is certified using its deployed expert records and the retained concavity/exit theorem. The controller has the explicit augmented state. Generation/checking are totals across nine genuinely distinct experts; initial checks and preprocessing remain in the ledger.\end{minipage}
\end{table}

In the analytic reference, class and temporal quadrature errors are zero, so the remaining gap is optimization error plus directed evaluation uncertainty. Direct L-BFGS-B reaches a worst gap below $3.393\times10^{-13}$ with a payoff enclosure width around $10^{-14}$. Neural Adam has one gap below $3.621\times10^{-8}$ and another near $.002772$, worse in the latter case than the direct first-order baselines. That adverse transferred-rate outcome is retained; it precludes universal neural superiority.
\begin{table}[htbp]\centering\small
\caption{Analytic reference: solver gap and arithmetic width}\label{tab:r24reference}
\begin{tabular}{lrrrrr}\toprule
Arm & Gap upper & Payoff width & Calls & Gen. s & Check s\\\midrule
Neural--Adam & 0.002771190 & $9.769963\times10^{-15}$ & 400.0 & 0.321 & 0.066\\
Direct--Adam & 0.00005492341 & $1.021406\times10^{-14}$ & 400.0 & 0.198 & 0.066\\
Diagonal--Adam & 0.00006450277 & $1.021406\times10^{-14}$ & 400.0 & 0.197 & 0.066\\
Tangent--Adam & 0.00001955686 & $1.021406\times10^{-14}$ & 400.0 & 0.197 & 0.067\\
Whitened--Adam & 0.5646302 & $1.332268\times10^{-14}$ & 400.0 & 0.196 & 0.069\\
Direct--L-BFGS-B & $3.392842\times10^{-13}$ & $1.021406\times10^{-14}$ & 14.0 & 0.007 & 0.067\\
\bottomrule\end{tabular}
\par\smallskip\begin{minipage}{.97\textwidth}\footnotesize Worst gap and width across two seeds, mean calls and times. This separate deterministic economy has zero class and temporal quadrature errors. Rates transfer from original-economy tuning without reference-specific calibration. The common consumption intercept is retained; the other two output channels are inactive.\end{minipage}
\end{table}


\subsection{Failure mechanism, verification slack and the unchanged objective}
The regenerated seed-23203 bound $.010503068$ reproduces the old failure. An equivalent layer-scaled chart reduces it to $.002554827$ within its 400-call diagnostic budget. Tighter tolerance and a direct restart from the already saturated raw outputs leave the bound near $.010503067$. Removing network redundancy only after reaching the adverse output is not sufficient; decoder saturation and the preceding trajectory also matter. Raw distances are not economic payoff distances. Restart budgets are explicitly additional, not equal-total-work comparisons.
\begin{table}[htbp]\centering\small
\caption{Post-hoc seed-23203 failure diagnosis}\label{tab:r24failure}
\begin{tabular}{lrrrr}\toprule
Case & Regret upper & Calls & Raw distance & Financed distance\\\midrule
regenerated & 0.01050307 & 92 & 299.963 & 2.691\\
scaled & 0.002554827 & 400 & 186.416 & 7.409\\
restart tighter & 0.01050307 & 43 & 296.812 & 2.753\\
direct restart & 0.01050307 & 4 & 299.963 & 2.691\\
\bottomrule\end{tabular}
\par\smallskip\begin{minipage}{.97\textwidth}\footnotesize Distances are to the unchanged frozen R23 direct solution. Tighter and direct restarts use additional budgets. Regenerated and chart-scaled trajectories start at identical outputs; these diagnostics are excluded from tuning and held-out selection.\end{minipage}
\end{table}

The largest held-out policy enclosure width is \RXXIVWidth{}. For exact policy payoff $J$ and unrestricted upper $\overline V$,
\[
 \overline V-L=(\overline V-V)+(V-J)+(J-L),\qquad 0\leq J-L\leq U-L.
\]
The first two terms are not separately identified in the original economy. Tight policy intervals alone do not identify dual slack. The reference separates these components in its own economy, while the original-economy pairwise intervals eliminate the dual. Preprocessing, generation, initial and checkpoint checking, scientific instrumentation, tuning and upper-library allocations are separately reported. The historical $415.789$-second dual cost is an accounting scenario, not a same-machine comparison; the reference/failure recovery stages also used a separately provisioned host.

The original-state all-start-time bound remains $7.241462443133396$ against the unchanged $.01$ objective. A new versioned record corrects the stale name of the older undiscounted field without changing historical bytes. This revision does not establish a new original-state full-domain accuracy certificate, a separate payoff increase for that actor, a successful nonreplicable stochastic execution, or rigorous numerical stopped-gradient errors. The new analytical and empirical results specify their actual domains while retaining the stronger objectives and their historical derivations.

